\pvs~\cite{pvs} is a verification system,
combining language expressiveness with automated tools. Its language is
based on higher-order logic, and is strongly typed.  It includes
types and terms for concepts such as: numbers, records, tuples, functions,
quantifiers, and recursive definitions.  Full predicate subtypes are
supported, which makes type checking undecidable; \pvs generates type
obligations (TCCs) as artefacts of type checking.  For example, division is
defined such that the second argument is nonzero, where nonzero is defined:
\begin{code}
nonzero_real: TYPE = {r: real | r /= 0}
\end{code}
Note that functions in \pvs are total; partiality is only supported via subtyping.

Beyond this, the \pvs language has structural subtypes (i.e., a record that
adds new fields to a given record), dependent types for records, tuples, and
functions, recursive and co-recursive datatypes, inductive and co-inductive
definitions, theory interpretations, and theories as parameters,
conversions, and judgements that provide control over the generation of
proof obligations.  Specifications are given as collections of
parameterized theories, which consist of declarations and formulas, and
are organized by means of imports. \medskip

The \pvs prover is interactive, but with a large amount of automation built
in.  It is closely integrated with the type checker, and features a
combination of decision procedures, including BDDs, automatic simplification,
rewriting, and induction.  There are also rules for ground evaluation,
random test case generation, model checking, and predicate abstraction.
The prover may be extended with user-defined proof strategies.

\pvs has been used as a platform for integration.  It has a rich API, making it relatively
easy to add new proof rules and integrate with other systems.  Examples of this include
the model checker, Duration Calculus, MONA, Maple, Ag, and Yices.  The system is normally
used through a customized Emacs interface, though it is possible to run it standalone
(PVSio does this), and \pvs features an XML-RPC server (developed independently of
the work presented here) that will allow for more flexible interactions.  \pvs is open source, and is
available at \url{http://pvs.csl.sri.com}.\medskip

As an example, \autoref{fig:intro:prel:pvs:equivclos} gives a part of the \pvs theory defining \emph{equivalence closures} on a type \expr{T} in its original syntax.
\pvs uses upper case for keywords and logical primitives; square brackets are used for types and round brackets for term arguments.
The most important declarations in theories are 
\begin{itemize}
	\item includes of other theories, e.g., the binary subset predicate \expr{subset?} and the type \expr{equivalence} of equivalence relations on \expr{T} are included from the theories \expr{sets} and \expr{relations} (These includes are redundant in the \pvs prelude and added here for clarity.), 
	\item typed identifiers, possibly with definitions such as \expr{EquivClos}, and 
	\item named theorems (here with omitted proof) such as \expr{EquivClosSuperset}.
\end{itemize}

\expr{VAR} declarations are one of several non-logical declarations: they only declare variable types, which can then be omitted later on; here \expr{PRED[[T,T]]} abbreviates the type of binary relations on \expr{T}.

\begin{figure}[ht]\centering
\begin{code}
EquivalenceClosure[T : TYPE] : THEORY
BEGIN
  IMPORTING sets, relations
  R: VAR PRED[[T, T]]
  x, y : VAR T
  EquivClos(R) : equivalence[T] =
    { (x, y) | FORALL(S : equivalence[T]) : subset?(R, S) IMPLIES S(x, y) }
  EquivClosSuperset : LEMMA
    subset?(R, EquivClos(R))
  ...
END EquivalenceClosure
\end{code}
\caption{An Excerpt From the \pvs Prelude Library}\label{fig:intro:prel:pvs:equivclos}
\end{figure}

%  EquivClosMonotone : LEMMA
%    subset?(R, S) IMPLIES subset?(EquivClos(R), EquivClos(S))


%%% Local Variables: 
%%% mode: latex
%%% TeX-master: "paper"
%%% End: 

%  LocalWords:  pvs subtyping parameterized organized Yices lstlisting emph centering
%  LocalWords:  fig:pvsequivclos EquivClos EquivClosSuperset mathhub
